\section{上游 ARA 原型的历史形式化}
\label{sec:method}

本节仅分析固定于 \code{edc3b123456c} 的上游原型。下述硬 $k$NN、全矩阵行归一化、五次 L-BFGS 调用、低秩路径仅重置 $B$ 与分段外层评分均为历史实现事实，不描述第~\ref{sec:ara-v1-results} 节的本地实验。本地 point-v1 的不同公式与控制流见第~\ref{sec:local-adaptation} 节。

\subsection{双层流程与问题定义}

本节把源实现还原为一个“冻结表示、内层改写、外层选择”的双层过程，而不把代码行为扩展为未经验证的训练算法。给定一组需要保持正常响应的提示 $\mathcal P_g$ 和一组需要降低拒答倾向的目标提示 $\mathcal P_b$，方法先在原模型上缓存所选模块的输入与输出，再独立修改各模块的线性映射，最后以模型级代理指标选择层区间和损失超参数。图~\ref{fig:ara-pipeline} 给出数据流：提示只在缓存阶段经过原模型，内层 L-BFGS 随后反复读取固定张量，不在每个闭包中重新运行完整 Transformer。该分离使参数更新可在模块尺度解释，也意味着内层目标只描述缓存点附近的行为，不能自动推出对任意提示的全局保证。源实现及其固定修订版本是本节历史算法事实的依据~\cite{ara2026source}。

设模型共有 $L$ 层，外层给出半开层区间 $[\ell_s,\ell_e)$，并在其中选择注意力输出投影和 MLP 下投影。对任一目标线性模块 $m$，记输入维为 $p$、输出维为 $q$，原权重为 $W_0\in\mathbb R^{q\times p}$。本节省略模块下标；实际过程按模块顺序重复。外层变量写为
\begin{equation}
  \theta=(\ell_s,\ell_e,\lambda_g,\lambda_b,\beta,k),
\end{equation}
其中 $\lambda_g$ 控制正常行为保持，$\lambda_b$ 控制目标提示的局部重映射，$\beta$ 控制远离原目标输出邻域的校正幅度。区间端点与连续参数的实际搜索域列于表~\ref{tab:source-parameters}；这些数值是源码默认值和搜索空间，不是本文重新调参所得的实验设置。

一次候选配置的执行顺序可以进一步展开为四步。首先，从未修改的基线模型收集所选层中全部目标模块的末词元记录，并按提示类别形成两个缓存字典。其次，对字典中的每个模块建立自己的优化变量，使用同一组 $\theta$ 最小化模块级目标；一个模块的损失不读取其他模块的权重，也不在上一模块改写后重新采集其输入。再次，将所有模块改写共同装入模型，在独立评估提示上计算外层两个目标。最后，外层搜索器记录该点并提出下一候选，模型快速恢复到基线映射后进入下一试次。于是，模块层面的拟合在冻结样本上可分解，模型层面的效果却由多个改写的组合决定；外层评分承担了检查这种组合结果的职责。

这一流程还区分了三种容易混淆的“方向”。参数优化当然产生梯度方向，L-BFGS 也保存历史搜索方向，但方法没有在表示空间中预先定义一个对所有目标样本共享的语义消融向量。真正决定每个目标输出移动参照的是它在 $Y_g$ 与 $Y_b$ 中的当前近邻，参照集合可因样本、模块和迭代而异。本文所称“无显式方向”专指没有固定的全局表示方向，而不是声称数值优化不使用方向信息，也不是声称任意修改都与子空间结构无关。

\begin{figure*}[t]
  \centering
  \resizebox{\textwidth}{!}{%
  \begin{tikzpicture}[
    font=\small,
    box/.style={draw, rounded corners=2pt, align=center, minimum height=11mm,
      minimum width=25mm, fill=blue!5, inner sep=4pt},
    cache/.style={box, fill=cyan!8},
    opt/.style={box, fill=orange!10},
    map/.style={box, fill=green!9},
    score/.style={box, fill=purple!8},
    arr/.style={-{Latex[length=2.2mm]}, thick},
    feedback/.style={-{Latex[length=2.2mm]}, thick, dashed, color=purple!70!black}
  ]
    \node[box] (prompt) {保持提示集 $\mathcal P_g$\\目标提示集 $\mathcal P_b$};
    \node[box, right=7mm of prompt] (hook) {目标层前向钩子\\仅取末词元};
    \node[cache, right=7mm of hook] (iocache) {冻结 I/O 缓存\\$X_g,Y_g,X_b,Y_b$};
    \node[opt, right=7mm of iocache] (loss) {$k$ 近邻三项损失\\$\mathcal L_{\rm total}$};
    \node[opt, right=7mm of loss] (lbfgs) {L-BFGS 内层优化\\强 Wolfe 线搜索};

    \node[map, above right=3mm and 8mm of lbfgs] (full) {全矩阵更新\\逐行范数保持};
    \node[map, below right=3mm and 8mm of lbfgs] (lora) {低秩更新\\$W_0+BA$};
    \node[score, right=12mm of $(full)!0.5!(lora)$] (score) {外层分段评分\\拒答比与首词元质量代理};
    \node[score, below=9mm of loss] (tpe) {多目标 TPE 提议 $\theta$\\更新层区间与损失超参数};

    \draw[arr] (prompt) -- (hook);
    \draw[arr] (hook) -- (iocache);
    \draw[arr] (iocache) -- (loss);
    \draw[arr] (loss) -- (lbfgs);
    \draw[arr] (lbfgs) -- (full);
    \draw[arr] (lbfgs) -- (lora);
    \draw[arr] (full) -- (score);
    \draw[arr] (lora) -- (score);
    \draw[feedback] (score.south) |- (tpe.east);
    \draw[feedback] (tpe.north) -- node[midway, right, align=left, font=\scriptsize]
      {下一试次：恢复基线映射\\复用冻结缓存} (loss.south);
  \end{tikzpicture}%
  }
  \caption{上游 ARA 原型的历史双层流程（硬近邻与分段评分）。前向钩子与 I/O 捕获只执行一次，外层试次恢复基线映射后复用冻结缓存；全矩阵与低秩路径是两种互斥的参数化方式。本图不描述本地 point-v1 的软目标、完整因子重置或原始指标评分。}
  \label{fig:ara-pipeline}
\end{figure*}


\subsection{冻结的模块输入输出}

缓存阶段利用前向钩子记录每个目标模块的成对输入输出。源实现对每条提示调用一次只生成一个新词元的前向过程，并从模块张量中仅取最后一个序列位置；记录值随即从计算图分离、复制并移至 CPU。对正常提示聚合后得到
\begin{equation}
  X_g\in\mathbb R^{n_g\times p},\qquad
  Y_g\in\mathbb R^{n_g\times q},
\end{equation}
对目标提示则得到 $X_b\in\mathbb R^{n_b\times p}$ 与 $Y_b\in\mathbb R^{n_b\times q}$。这里 $X$ 是进入线性层的末词元隐藏状态，$Y$ 是同一次原始前向中的模块输出；四个张量在相应模块被优化期间保持冻结。因此，后续预测 $XW^{\mathsf T}$ 是针对原始模块局部输入分布的离线回放，而不是随权重变化重新采样的端到端轨迹。

冻结机制还带来两个必须显式满足的前置条件。第一，批量聚合允许某些钩子在部分批次中没有出现，但两类缓存随后按相同模块键直接索引，故被优化模块必须同时存在于正常缓存和目标缓存。第二，两个近邻集合分别从 $Y_g$ 和 $Y_b$ 中选取 $k$ 个样本，所以至少要求 $n_g\geq k$ 且 $n_b\geq k$。若任一条件不成立，当前实现没有定义可替代的跨模块匹配或自动缩小 $k$ 的规则。把这些条件写入形式化描述，可以区分算法所需不变量与通常数据批次恰好满足的偶然情况。

末词元采样决定了缓存所覆盖的条件分布。每条提示在生成调用中只允许产生一个新词元，钩子仍可能在预填充与生成相关计算中被触发，但保存操作统一截取当前张量的最后位置，并按样本维拼接。因而 $X_g$ 与 $X_b$ 不是整段序列所有位置的集合，$Y_g$ 与 $Y_b$ 也不是完整生成轨迹。CPU 分离副本阻断了回到基线模型的梯度，优化图只包含待改写权重、线性重放、距离计算与损失。这个设计降低了闭包成本，同时把结论范围限定在末词元模块响应：若要研究长序列位置间的传播，必须另行采样，而不能从当前缓存直接推得。

缓存中的输入输出也具有不同角色。保持项使用同一正常样本的 $(x_{g,i},y_{g,i})$ 配对，保留了逐样本监督；两个距离项只把输出行视为参考点集，不要求 $x_{b,i}$ 与某个正常提示一一配对。由此，方法可以在两类提示数量不同的情况下工作，但不能把近邻索引解释为提示语义对应关系。尤其是多个目标查询可以选择同一个正常输出，参考点也不需要被均匀使用，这正是其与带容量约束匹配方法的关键区别。

\subsection{局部近邻三项目标}

局部目标以动态 $k$ 近邻距离代替预先估计的单一方向。对查询向量 $z\in\mathbb R^q$ 和参考矩阵 $Y=(y_1,\ldots,y_n)^{\mathsf T}$，定义欧氏距离下的近邻索引集合及平均距离为
\begin{align}
  \mathcal N_k(z;Y)
    &= \operatorname*{arg\,topk}_{j\in\{1,\ldots,n\}}^{\rm smallest}
       \lVert z-y_j\rVert_2, \\
  d_k(z;Y)
    &= \frac{1}{k}\sum_{j\in\mathcal N_k(z;Y)}
       \lVert z-y_j\rVert_2 .
  \label{eq:knn-distance}
\end{align}
源实现通过全对全的 \code{cdist} 生成距离矩阵，再以最小值 \code{topk} 选择邻居。邻居索引随当前权重产生的查询位置重新计算，不是缓存阶段固定的标签。

正常保持项要求新映射在正常缓存上复现原模块输出：
\begin{equation}
  \mathcal L_{\rm good}(\widetilde W)
  =\frac{1}{n_gq}\left\lVert
    X_g\widetilde W^{\mathsf T}-Y_g
  \right\rVert_F^2 .
  \label{eq:good-loss}
\end{equation}
目标提示的输出一方面被拉向正常输出的局部邻域，另一方面被推离其原目标输出邻域：
\begin{align}
  \mathcal L_{\rm pull}(\widetilde W)
    &=\frac{1}{n_b}\sum_{i=1}^{n_b}
      d_k(x_{b,i}\widetilde W^{\mathsf T};Y_g), \\
  \mathcal L_{\rm push}(\widetilde W)
    &=\frac{1}{n_b}\sum_{i=1}^{n_b}
      d_k(x_{b,i}\widetilde W^{\mathsf T};Y_b).
  \label{eq:pull-push}
\end{align}
三项按
\begin{equation}
  \begin{aligned}
  \mathcal L_{\rm total}(\widetilde W)
    ={}&\lambda_g\mathcal L_{\rm good}(\widetilde W)\\
       &+\lambda_b\bigl[
         \mathcal L_{\rm pull}(\widetilde W)
         -\beta\mathcal L_{\rm push}(\widetilde W)
       \bigr]
  \end{aligned}
  \label{eq:total-loss}
\end{equation}
组合。负号使最小化过程增大到 $Y_b$ 的邻域距离；$\mathcal L_{\rm good}$ 则限制这一变化对正常缓存映射的偏离。该构造不先求“拒答方向”，而让每个查询使用邻近正常输出形成局部参照，因而可覆盖多簇表示而无需把它们压缩为单个均值差向量。

各超参数的作用可以直接由式~\eqref{eq:total-loss} 判读。$\lambda_b=0$ 时，目标提示不再参与改写；$\beta=0$ 时，目标输出只被拉向正常邻域而不显式远离原邻域；增大 $k$ 会把单点参照换成更宽的局部平均，但并不保证距离更平滑，因为近邻集合仍可切换。$\mathcal L_{\rm good}$ 是按元素平均的平方距离，$\mathcal L_{\rm pull}$ 与 $\mathcal L_{\rm push}$ 是按样本平均的一阶欧氏距离，三者量纲和随 $q$ 的尺度并不相同。因此，$\lambda_g$、$\lambda_b$ 与 $\beta$ 不应被解读为概率或归一化混合系数；外层搜索实际承担了在这些不同尺度间寻找折中的任务。

局部参照比全局中心保留了更多多峰结构，但它也会继承缓存密度的不均衡。高密度正常簇可能成为更多查询的近邻，孤立点的影响则受 $k$ 和相对距离限制；推远项同样优先作用于当前最接近的原目标输出，而非整个目标分布。方法由此表达的是经验点云上的局部几何，不是两个真实表示分布之间的总体距离。若缓存样本不足或分布偏置，目标函数本身不会校正这种偏置，外层有限评估也只能提供代理层面的筛选。

这里的“局部邻域映射”不是最优传输。式~\eqref{eq:knn-distance} 没有样本间耦合矩阵、质量守恒约束或 Wasserstein 代价，只是逐查询的近邻选择；即使采用“局部传输视角”，也只能作为几何解释，不能据此声称求解了最优传输问题。动态邻域还使全局目标非光滑且通常非凸。例如 $k=1$、$Y=\{-1,1\}$ 时，$d_1(z;Y)=\min(|z+1|,|z-1|)$ 在两个谷底的中点违反凸性不等式。邻居切换边界和零距离处也可能不可微。因此，后文的拟牛顿过程应被理解为源码采用的数值求解器，而不是全局最优性证明。

\subsection{全矩阵更新与逐行归一化}

全矩阵路径把 $W\in\mathbb R^{q\times p}$ 本身设为待优化变量，因此没有显式低秩约束；“全矩阵”只描述参数化自由度，不保证最终改变量具有满秩。为限制各输出通道的尺度，源实现先保存初始行范数 $s_j=\lVert w_{0,j}\rVert_2$，并在每次闭包计算中使用重参数化代理
\begin{equation}
  \mathcal R(W)_{j,:}
  =s_j\frac{w_{j,:}}
  {\max(\lVert w_{j,:}\rVert_2,\varepsilon)},
  \qquad \widetilde W=\mathcal R(W).
  \label{eq:row-normalization}
\end{equation}
在 $\lVert w_{j,:}\rVert_2\geq\varepsilon$ 的常规区域，代理行范数严格等于 $s_j$；在数值保护分支中，其范数不大于 $s_j$，当 $s_j>0$ 时严格小于。因此不能把式~\eqref{eq:row-normalization} 无条件描述成等式约束投影，也不宜把可行域概括为一个整体球面。内层计算 $X\widetilde W^{\mathsf T}$，没有在重建公式中另加显式偏置；优化结束后，再把归一化后的矩阵复制回模块权重。

逐行处理保留的是初始权重每个输出通道的尺度，而不是整个矩阵的单一 Frobenius 范数。常规区域中的自由度主要来自行方向旋转；不同输出行可以独立改变方向，行与行之间没有正交约束。若初始某行范数为零，则对应 $s_j=0$，代理行持续为零；若待优化行落入保护区域，分母被截断，梯度行为也不同于单位球面上的光滑重参数化。与此同时，钩子缓存的 $Y$ 是模块真实输出，而重放公式只写 $X\widetilde W^{\mathsf T}$，故可能存在的固定偏置并未作为单独变量加入目标。保持项会把这部分差异一并吸收到权重拟合压力中，不能说它精确隔离了权重变化之外的所有因素。

每个模块使用有限记忆 BFGS 及强 Wolfe 线搜索~\cite{liu1989lbfgs}。源码设置学习率为 $1$、每次 \code{step} 的内部最大迭代数为 $20$、历史长度为 $10$，并固定执行五次 \code{step(closure)}。一次闭包会重新构造归一化代理、全对全距离和近邻索引，然后反向传播式~\eqref{eq:total-loss}。固定调用次数只是实现预算；本文没有迭代轨迹或最优性证书，因而不把“五次”解释为已收敛。

五次外部调用也不能简单换算为一百次闭包。$20$ 是每次调用允许的内部迭代上限，强 Wolfe 搜索可为试探步额外求值，而提前停止也可能减少求值；实际次数应以 $C_m$ 记账。L-BFGS 的曲率近似通常适合中等规模平滑问题，但此处近邻切换和欧氏范数零点引入非光滑性，逐行保护分支又改变局部导数。采用该求解器是可复现的工程选择，不足以建立标准光滑拟牛顿理论中的局部收敛前提。

\subsection{低秩工程变体}

低秩路径冻结 $W_0$，以
\begin{equation}
  W_{\rm eff}=W_0+BA,\qquad
  B\in\mathbb R^{q\times r},\quad
  A\in\mathbb R^{r\times p}
  \label{eq:lora}
\end{equation}
表示有效权重，其增量秩至多为 $r$~\cite{hu2022lora}。源默认秩为 $128$，且设置 $\alpha=r$，所以常见 LoRA 写法中的缩放 $\alpha/r$ 等于 $1$。模块配置中的无偏置选项表示不训练适配器偏置参数，不能据此推断原线性模块一定没有固有偏置。

该变体存在需要审计的目标与部署差异。当行归一化开关开启时，闭包把 $\mathcal R(W_0+BA)$ 送入损失；部署时适配器实际给出的却是未显式归一化的 $W_0+BA$。两者不相等时，优化代理与最终映射并非同一函数。另一个状态条件是快速重置只把 $B$ 清零而保留 $A$：有效映射由此恢复到 $W_0$，但后续非凸优化的参数初始化可能依赖先前试次。本文保留这些源码语义，不把低秩路径与全矩阵路径表述成完全等价的两个求解器。

从表达能力看，全矩阵路径允许 $qp$ 个权重元素共同变化，低秩路径只有 $r(p+q)$ 个适配器元素，并把改变量限制在秩不超过 $r$ 的集合中。参数更少不等于有效映射无需完整矩阵：闭包仍先形成 $W_0+BA$，再进行与全矩阵相同的批量乘法和邻域计算。初始化时只要 $B=0$，有效映射便等于 $W_0$，但 $A$ 决定从该函数点出发的参数化切空间；保留 $A$ 因而可能改变下一试次早期梯度。比较两个路径时，应同时报告参数预算、归一化是否开启以及试次重置策略，而不能只报告秩。

\subsection{外层分段评分与计算代价}

外层以多变量 TPE 在表~\ref{tab:source-parameters} 的空间中产生 $\theta$，并同时最小化质量代理与拒答代理~\cite{akiba2019optuna}。设基线模型在评估目标提示上的拒答计数为 $f_0$，修改后计数为 $f$，则归一化拒答量为
\begin{equation}
  r_f=\begin{cases}
    f/f_0, & f_0>0,\\
    f, & f_0=0.
  \end{cases}
  \label{eq:refusal-ratio}
\end{equation}
质量变化使用首个生成词元分布的
$D=D_{\rm KL}(p_{\rm base}\Vert p_{\rm edit})$。源码传给搜索器的两个目标为
\begin{equation}
  s_1=\begin{cases}
    D/\sigma, & D\geq\tau,\\
    r_f\tau/\sigma, & D<\tau,
  \end{cases}
  \qquad s_2=r_f,
  \label{eq:outer-score}
\end{equation}
默认 $\sigma=1$、$\tau=0.01$。低于 KL 阈值时，第一目标被拒答量替换为按拒答代理缩放的惩罚，而不是继续直接使用 $D$。可选的通用能力分支以负的归一化准确率替代 $s_1$。还应区分搜索目标与结果展示：后者保存原始拒答计数和原始 KL（或负准确率），不等同于式~\eqref{eq:outer-score} 的分段坐标。这些指标只是拒答行为与局部质量的代理，不能单独证明全面安全性或通用能力保持。

分段设计把阈值内候选的首要区分依据改为拒答量：当 $D<\tau$ 时，$s_1$ 与 $s_2$ 都随 $r_f$ 变化，只是尺度不同；到达阈值后，$s_1$ 才恢复为 KL。它不是 $D$ 与 $r_f$ 的固定加权和，在分支边界也不必连续。例如 $r_f<1$ 时，从阈值下方接近所得 $r_f\tau/\sigma$ 小于阈值处的 $\tau/\sigma$。因此，Pareto 支配关系由这条显式策略塑造，读取搜索结果时必须保留分支条件。基线计数 $f_0$ 只评估一次，零基线使用原始 $f$ 的退化分支避免除零；这也使不同数据切分或不同基线之间的数值尺度未必可直接比较。

首词元 KL 的方向同样需要固定：实现计算的是 $D_{\rm KL}(p_{\rm base}\Vert p_{\rm edit})$，不是反向 KL，也不是整段生成序列的散度。它衡量基线分布质量在修改分布下的相对对数概率变化，对修改分布在基线低概率区域增加的质量敏感性有限。拒答计数则依赖外部判别规则，把连续输出压缩为离散事件。二者结合能为外层搜索提供便宜的二维信号，但没有覆盖事实正确性、多轮一致性、鲁棒性或广泛任务能力；可选准确率分支也只增加一个特定任务代理。

\begin{table*}[t]
  \centering
  \caption{上游 ARA 原型的历史默认值与搜索域（edc3b123456c）。所有行仅属于该原型，不是本地 point-v1 或 trajectory-v2 的实验配置。}
  \label{tab:source-parameters}
  \footnotesize
  \setlength{\tabcolsep}{4pt}
  \renewcommand{\arraystretch}{1.12}
  \begin{tabularx}{\textwidth}{@{}p{22mm}p{32mm}p{58mm}Y@{}}
    \toprule
    类别 & 参数 & 默认值或搜索域 & 实现语义 \\
    \midrule
    目标模块 & 组件 & 注意力输出投影、MLP 下投影 & 在所选连续层区间内逐模块建立缓存并优化 \\
    层区间 & 起始层、结束层 & $[0,\lfloor L/2\rfloor]$、$[\lfloor L/2\rfloor,L]$ & 结束层按 Python 切片语义不包含在区间内 \\
    损失权重 & $\lambda_g,\lambda_b,\beta$ & $[0,1]$；$[10^{-4},1]$（对数采样）；$[0,1.3]$ & 分别控制保持项、目标项和远离原目标邻域的校正强度 \\
    局部邻域 & $k$ & 整数 $[1,15]$ & 两个距离项使用相同近邻数 \\
    全矩阵路径 & 行归一化 & 默认开启 & 每次闭包计算前按初始权重的逐行范数重参数化 \\
    低秩路径 & 开关、秩 $r$ & 默认关闭；$r=128$ & 缩放取 $\alpha=r$，故适配器乘子 $\alpha/r=1$ \\
    L-BFGS & 学习率、内部迭代、历史、线搜索 & $1$；$20$；$10$；strong Wolfe & 每个目标模块固定调用 $5$ 次 \code{step(closure)} \\
    外层搜索 & 试次数、启动试次、候选数 & $200$；$60$；$128$ & 多变量 TPE，两个目标均最小化 \\
    质量代理 & $\sigma,\tau$ & $1.0,0.01$ & 用于首词元 KL 的缩放与分段阈值 \\
    缓存规模 & 优化/评估提示数 & 每类 $400$；每类 $100$ & 前者建立内层缓存，后者计算外层代理 \\
    \bottomrule
  \end{tabularx}
\end{table*}


计算瓶颈来自代理矩阵乘法和全对全近邻。对单模块单闭包，前者时间为 $O((n_g+n_b)pq)$；拉近与推远距离分别需要 $O(n_bn_gq)$ 和 $O(n_b^2q)$ 时间，并分别物化 $O(n_bn_g)$ 与 $O(n_b^2)$ 的距离矩阵。由于不同后端的 \code{topk} 实现不同，选择开销保守记为
$T_{\rm topk}(n_b,n_g,k)+T_{\rm topk}(n_b,n_b,k)$，而不强行指定排序复杂度。若模块 $m$ 的实际闭包调用数为 $C_m$，则主导总时间可写为
\begin{equation}
  O\!\left(\sum_m C_m\left[
    (n_g+n_b)p_mq_m+n_b(n_g+n_b)q_m+T_{{\rm topk},m}
  \right]\right).
\end{equation}
全矩阵 L-BFGS 的历史状态量为 $O(hp_mq_m)$，其中 $h=10$；低秩可训练状态降为 $O(hr(p_m+q_m))$，但仍需物化 $W_0+BA$ 的 $O(p_mq_m)$ 空间，并承担形成 $BA$ 的 $O(p_mq_mr)$ 时间。CPU 缓存总标量数为
$O(\sum_m(n_{g,m}+n_{b,m})(p_m+q_m))$。因此低秩路径主要压缩优化变量和拟牛顿历史，而不会消除完整有效权重、缓存或距离矩阵的代价。以上分析完全来自固定源码和静态重建；本节未新增模型运行或经验实验。

外层总成本还要乘以被实际评估的候选配置，并随候选层区间内的模块数变化。默认两百试次是搜索预算而非保证完成的有效样本数；失败、剪枝或运行环境差异都可能改变实际求值集合。距离矩阵的峰值显存随 $n_b(n_g+n_b)$ 二次或双线性增长，低秩化并不改变这一项，因而在大缓存上可能先受邻域内存限制而不是参数历史限制。相反，缓存常驻 CPU，只有当前模块张量参与设备侧闭包，减轻了同时保存所有模块激活的设备压力。完整资源报告至少应给出每类样本数、目标模块数、实际闭包次数和外层有效试次数；源实现中的默认值只能用于复现配置，不能替代实测时间与峰值内存。
